\documentclass[12pt]{article}
\usepackage[utf8]{inputenc}
\usepackage[T1]{fontenc}
\usepackage[margin=1in]{geometry}
\usepackage{setspace}
\usepackage{xcolor}
\usepackage[colorlinks=true, urlcolor=blue!60!black, linkcolor=blue!60!black]{hyperref}
\setlength{\parindent}{0em}
\setlength{\parskip}{0.45em}
\renewcommand{\baselinestretch}{1.0}
\begin{document}
\begin{center}
{\Large \textit{Korematsu v. United States}, 323 U.S. 214 (1944)\par}
\vspace{0.3em}
{\normalsize Case Brief --- One-Page Edition\par}
\end{center}

\textbf{Memory Jogger.} The Japanese American exclusion case: wartime deference despite ``immediately suspect'' racial classifications, now repudiated (\href{https://www.courtlistener.com/opinion/104040/korematsu-v-united-states/}{323 U.S. 214}).

\textbf{Facts.} Under Executive Order 9066, a military order excluded all persons of Japanese ancestry, including citizens, from part of California. Fred Korematsu, a loyal citizen, stayed home and was arrested.

\textbf{Procedural History.} Korematsu was convicted and placed on probation; the Ninth Circuit affirmed. The conviction was vacated on \textit{coram nobis} in 1984.

\textbf{Votes.} Affirmed, 6--3. Black wrote for the Court; Frankfurter concurred; Roberts, Murphy, and Jackson dissented.

\textbf{Holding.} Racial restrictions are ``immediately suspect,'' but when Congress and the President act jointly in wartime against grave danger, temporary exclusion of citizens of Japanese ancestry is permissible. The Court did not decide detention.

\textbf{Reasoning.} Black deferred to military judgment that the disloyal could not be quickly separated from the loyal and denied racial hostility. Roberts saw a ``cleverly devised trap''; Murphy found racial stereotypes, not evidence; Jackson warned that judicial approval creates a lasting principle. \textit{Trump v. Hawaii} later called the decision ``gravely wrong the day it was decided.''

\textbf{Examples.} \emph{Same-Side:} detaining citizens of one origin without hearings. \emph{Opposite-Side:} a statutory entry ban on foreign nationals. \emph{Fence-Sitter:} travel-based screening that falls on one ethnic group.

\textbf{Critique.} A fixture of the ``anticanon''; Harold Hongju Koh warns that national-security deference risks reviving its ``ghost.''

\textbf{Key Quotation.} ``It should be noted, to begin with, that all legal restrictions which curtail the civil rights of a single racial group are immediately suspect. That is not to say that all such restrictions are unconstitutional. It is to say that courts must subject them to the most rigid scrutiny. Pressing public necessity may sometimes justify the existence of such restrictions; racial antagonism never can.'' --- Justice Black, Opinion of the Court.

\end{document}
